\documentclass{beamer}

\mode<presentation>{
  \usetheme{Madrid}
  \usecolortheme{default}
  \usefonttheme{default}
  \setbeamertemplate{navigation symbols}{}
  \setbeamertemplate{caption}[numbered]
}

% ===== Idioma y codificación =====
\usepackage[utf8]{inputenc}
\usepackage[T1]{fontenc}
\usepackage[spanish]{babel}

% ===== Paquetes que SÍ usas =====
\usepackage{graphicx}    % imágenes
\usepackage{amsmath}     % matemáticas
\usepackage{circuitikz}  % diagrama eléctrico (circuito)
\usepackage{tikz}        % por si deseas otro dibujo simple (no abusamos)
\usetikzlibrary{positioning}

% ===== Comandos personalizados (y los usamos más abajo) =====
\newcommand{\materia}[1]{\textbf{\textcolor{blue}{#1}}}
\newcommand{\miNombre}{\textbf{Brandon Avalos}}

% ===== Datos portada =====
\title[Mecatrónica]{Mecatrónica}
\subtitle{Día de un estudiante}
\author[\miNombre]{\miNombre}
\institute{Facultad de Estudios Profesionales Zona Media}
\date{25/02/2026}

\begin{document}

% ---------- Frame 1: Portada ----------
\begin{frame}
  \titlepage
\end{frame}

% ---------- Frame 2: Tabla + Imagen ----------
\begin{frame}{Mi Rutina Semanal}
{\footnotesize
\centering
\begin{tabular}{|c|c|c|c|c|}
\hline
Lunes & Martes & Miércoles & Jueves & Viernes \\ \hline
Espalda/Bíceps & Pierna completa & Pecho/Tríceps & Espalda/Bíceps & Pierna completa \\ \hline
\end{tabular}
\par\medskip
\includegraphics[width=0.45\linewidth]{espal.png} % <-- renombra o sube tu imagen con este nombre
\captionof{ }{  }
}
\end{frame}

% ---------- Frame 3: Introducción (usa comandos personalizados) ----------
\begin{frame}{Introducción}
Presento un resumen de mi día como estudiante de \materia{Mecatrónica}.  
Autor: \miNombre.
\end{frame}

% ---------- Frame 4: Tabla simple adicional (clarita) ----------
\begin{frame}{Material básico}
\centering
\begin{tabular}{l|r}
objeto & Cantidad \\ \hline
Cuaderno & 2 \\
Pluma & 3 \\
Calculadora & 1 \\
\end{tabular}
\captionof{ }{ }
\end{frame}

% ---------- Frame 5: Matemáticas (mínimo contenido académico) ----------
\begin{frame}{Ejemplo Matemático}
Sea $X_1, X_2,\ldots,X_n$ i.i.d. con media $\mu$ y varianza $\sigma^2 < \infty$.
\[
S_n = \frac{1}{n}\sum_{i=1}^{n} X_i
\]
Entonces:
\[
\sqrt{n}(S_n - \mu) \xrightarrow{d} \mathcal{N}(0,\sigma^2)
\]
\end{frame}

% ---------- Frame 6: Diagrama/Dibujo (Circuito eléctrico simple) ----------
\begin{frame}{Diagrama eléctrico sencillo}
\centering
\begin{circuitikz}
\draw
(0,0) to[battery1,l=$V$] (0,2)
      to[R,l=$R$] (3,2)
      -- (3,0) -- (0,0);
\end{circuitikz}
\end{frame}

\end{document}